\section{Summary}
\label{sec:comp-asymptotic-summary}
Use $X=(s_{\mathrm{phys}}-m^2)/m^2=2\omega/m$, where
$\omega$ is the incident photon energy in the electron rest frame.
Write the NLO total as $\sigma_0+\sigma_1$, with $\sigma_0$ of
order $\alpha_{\mathrm T}^2$ and $\sigma_1$ of order
$\alpha_{\mathrm T}^3$. The low-energy results are
\begin{equation}
 \sigma_0=\frac{\pi\alpha_{\mathrm T}^2}{m^2}
 \left[\frac83-\frac83X+\frac{52}{15}X^2
                   -\frac{133}{30}X^3+O(X^4)\right].
\end{equation}
\begin{equation}
\begin{aligned}
 \sigma_1&=\frac{\alpha_{\mathrm T}^3}{m^2}
 \biggl\{X^2\left[-\frac{16}{9}\ln X+\frac7{15}\right]
 +X^3\left[\frac43\ln X-\frac{113}{90}\right]\\
 &\hspace{10mm}+X^4\left[-\frac{25}{9}\ln X+\frac{203}{125}\right]
                   +O\!\left(X^5|\ln X|\right)\biggr\},
 \qquad X\to0^+.
\end{aligned}
\end{equation}
The combination of virtual, real, and soft contributions is derived in
Section~\ref{sec:comp-independent-total-low}, and the further orders follow
directly from the collected formula in Section~\ref{sec:comp-total-energy-low}.
The photon-mass logarithms
cancel before taking this limit. The radiative correction vanishes,
while the Born cross section tends to $8\pi\alpha_{\mathrm T}^2/(3m^2)$.

At high energy, the Born result is
\begin{equation}
 \sigma_0=\frac{\pi\alpha_{\mathrm T}^2}{m^2}
 \left[\frac{2\ln X+1}{X}+\frac{18-8\ln X}{X^2}
                +O\!\left(\frac{\ln X}{X^3}\right)\right],
 \qquad X\to\infty.
\end{equation}
The NLO correction is
\begin{equation}
\begin{aligned}
 \sigma_1=\frac{\alpha_{\mathrm T}^3}{m^2}\biggl\{
 \frac1X\biggl[&\frac{\ln^3X}{3}-\frac{\ln^2X}{2}
 +\frac{17}{4}\ln X-\frac{75}{8}-\frac{\pi^2}{2}+4\zeta(3)\biggr]
 +O\!\left(\frac{\ln^3X}{X^2}\right)\biggr\}.
\end{aligned}
\end{equation}
Section~\ref{sec:comp-total-energy-high} derives this expansion directly
from the collected total, including the finite constants of its GPLs.
The comparison with the published expansion is made after those
constants are determined.
